\section{Differenzialgleichungen}

\begin{tcolorbox}
	Gegeben $F$, eine Funktion von $x,y$, und die Ableitung von $y$, wobei $y$ selbst eine Funktion ist.
	Dann ist die Gleichung der Form
	$$F(x,y,y', \dots, y^{(n)}) = b(x)$$
	eine \textbf{gewöhnliche Differenzialgleichung} (kurz DGL / ODE) von \textbf{Ordnung} $n$. Alternativ:
	$$y^{(n)} + a_{n-1}y^{(n-1)} + \dots + a_1 y' + a_0 y = b(x)$$
	wobei die Koeffizienten $a_0, \dots, a_{n-1}$ komplexwertige Funktionen auf $I \subset \R$ sind (können
	von $y$ abhängig sein). Wenn $b(x) = 0$, 
	dann ist das DGL \textbf{homogen}, sonst ist es inhomogen. 
\end{tcolorbox}

Eine Differenzialgleichung ist \textbf{linear} wenn $F$ als eine lineare Kombination von den Ableitungen von $y$ 
geschrieben werden kann:

$$b(x) = \Sum_{i = 0}^{n}a_i(x) \cdot y^{(i)}$$

wobei $a_i(x)$ und $b(x)$ stetige Funktionen sind.

Erkennen von linearen ODE:
\begin{enumerate}[label=(\arabic*)]
    \item keine Koeffizienten vor der höchsten Ableitung
    \item alle Koeffizienten sind stetige Funktionen
    \item keine Produkte von $y$ oder deren Ableitungen
    \item keine Potenzen von $y$ oder deren Ableitungen
    \item keine Funktionen die von $y$ oder deren Ableitungen abhängen
\end{enumerate}

Die Menge $S$ der Lösungen von $f$ zu einem DGL, ist eine Teilmenge des Raums der komplexwertigen Funktionen auf 
$I$, mit der Dimension $n$. Im Allgemeinen sind solche Lösungen nicht eindeutig, nur mit Anfangsbedingungen können
wir eine eindeutige Lösung erhalten. Ein DGL der Ordnung $k$ benötigt $k$ Anfangsbedingungen für eine eindeutige
Lösung. \medskip 

$S_0$ ist die \textbf{Menge der Lösungen zum homogenen DGL}. Für ein beliebiges $b$ ist die Menge 
der Lösungen gegeben durch $S_b = \{f_h + f_p \;|\; f_h \in S_0\}$ wobei $f_p$ eine 
partikulär Lösung ist. \medskip

\Title{Lineare DGL von Ordnung 1}

Betrachten wir lineare DGL von der Form:
$$y' + a(x)y = b(x)$$

\begin{enumerate}[label=(\arabic*)]
\item Zuerst lösen wir das homogenen DGL:
\begin{align*}
    y' + a(x)y &= 0\\
    \frac{y'}{y} &= -a(x)\\
    \ln(y) &= -A(x) \light{+C}\\
    y &= e^{-A(x) \light{+C}} = z \cdot e^{-A(x)}\quad z \in \C
 \end{align*}
Bemerke, dass wenn $f$ eine Lösung ist, muss auch $z \cdot f, \forall z \in \C$ eine Lösung sein.

\item Nun finden wir eine partikulär Lösung $f_p: I \to \C$ so dass $f_p' + a(x)f_p = b(x)$. Dazu verwenden wir "Variation
der Konstanten" oder "fundiertes Raten".
\end{enumerate}

Das Anfangswertproblem mit $f(x_0) = y_0$ hat eine eindeutige Lösung, gegeben durch $f(x) = y_0 \cdot e^{A(x_0) - A(x)}$.


\Title{Fundiertes Raten}

Wenn $b(x)$ von einer bestimmten Form ist, versuchen wir folgende $f_p$:

\begin{center}
  \begin{tabular}{|c|c|} 
   \hline
    $b(x)$ & Raten \\ \hline
    $a \cdot e^{\alpha x}$ & $b \cdot e^{\alpha x}$\\ \hline
    $a \sin(\beta x)$ or $b \cos(\beta x)$  & $c \sin(\beta x) + d \cos(\beta x)$\\ \hline
    $a e^{\alpha x} \sin(\beta x)$ or $b e^{\alpha x} \cos(\beta x)$ & $e^{\alpha x} \Big( c \sin(\beta x) + d \cos(\beta x) \Big)$\\ \hline
    $P_n(x) e^{\alpha x}$ & $R_n(x) \cdot e^{\alpha x}$\\ \hline
    $P_n(x) e^{\alpha x} \sin(\beta x)$ & $e^{\alpha x} \Big( R_n(x) \sin(\beta x) + S_n(x) \cos(\beta x) \Big)$\\ \hline
    $P_n(x) e^{\alpha x} \cos(\beta x)$ & $e^{\alpha x} \Big( R_n(x) \sin(\beta x) + S_n(x) \cos(\beta x) \Big)$\\ \hline
  \end{tabular}
\end{center}


\begin{enumerate}[label=(\arabic*)]
    \item Wenn $b(x)$ eine Linearkombination der Basisfunktionen ist, so versuche eine Linearkombination.
    \item Wenn die geratene Lösung die selbe ist wie die homogene Lösung, so versuche sie mit $x^m$ zu 
    multiplizieren, wobei $m$ die Vielfachheit der Wurzel ist.
\end{enumerate}


\Title{Variation der Konstanten}

\begin{enumerate}[label=(\arabic*)]
  \item Nimm an $f_p = z(x) e^{-A(x)}$ für eine Funktion $z: I \to \C$
  \item Wir setzen dies in die Gleichung ein und schauen was $z$ sein kann
  \begin{align*}
    y' + a(x)y &= b(x)\\
    \Rightarrow z'(x) e^{-A(x)} &= b(x)
  \end{align*}
  
  Daraus folgt, dass: 
  \begin{align*}
    z'(x) &= b(x) e^{A(x)}\\
    z(x) &= \int_{x_0}^x b(t)e^{A(t)}dt 
  \end{align*}
 
  Daher erhalten wir $f_p = \int_{x_0}^x b(t)e^{A(t)}dt  \cdot e^{-A(x)}$.
\end{enumerate}


\Title{Lineare DGL mit konstanten Koeffizienten}

Betrachten wir DGL von der Form:

$$y^{(n)} + a_{n-1}y^{(k-1)} + \dots + a_1y' + a_0y = b(x)$$

wobei $a_0, \dots, a_{n-1} \in \C$ konstant sind und $k \geq 1$. $b(x)$ ist weiterhin eine stetige Funktion.
Die Menge der Lösungen $S$ ist somit ein $\C$-Vektorraum mit dim($S) = k$. \medskip

Wir suchen zuerst die \textbf{homogenen Lösung}. Dafür nehmen wir an, dass die Lösung von der Form 
$e^{\lambda x}$ ist. Nun können wir das charakteristische Polynom lösen:
\begin{align*}
  P(\lambda) &= \lambda^k e^{\lambda x} + a_{k-1} \lambda^{k-1} e^{\lambda x} + \dots + a_0 e^{\lambda x} = 0\\
  \Rightarrow 0 &= \lambda^k + a_{k-1}\lambda^{k-1} + \dots + a_0
\end{align*}
Die Nullstellen von $P(\lambda)$ sind Eigenwerte, $\lambda_i$, mit dazugehöriger Multiplizität  $m_r$. Nun spannen
die Funktionen $f_{i,r} : x \to x^r e^{\lambda_i x}$ den Raum der Lösungen $S_0$ auf. \medskip

Falls $\lambda = \beta + \gamma i$ eine komplexe Nullstelle von $P(\lambda)$ ist, so ist auch 
$\bar{\lambda} = \beta - \gamma i$ eine Nullstelle. Daher sind $f_1 = e^{\lambda x}$ und
$f_2 = e^{\bar{\lambda} x}$ Lösungen und könen durch eine Linearkombination von 
$\tilde{f_1} = e^{\beta x} \cos(\gamma x)$ und $\tilde{f_2} = e^{\beta x} \sin(\gamma x)$.
ersetzt werden. \medskip

Falls $y^{(k)} + a_{k-1}y^{(k-1)} + \dots + a_0 y = 0$ reelle Koeffizienten hat, so führt jedes paar von komplex
konjugierten Nullstellen $\beta_j \pm \gamma_j i$ mit Multiplizität $m_j$ zu einer Lösung

$$ x^l e^{\beta_j x} \Big( \cos(\gamma_j x) + i \sin(\gamma_j x) \Big) \quad \text{für }0 \leq l \leq m_j$$

Um eine \textbf{partikulär Lösung} zu finden, können wir wieder "fundiertes Raten" oder "Variation der 
Konstanten" verwenden. Wir verwenden "Variation der Konstanten", wenn $k = 2$ und die RHS "unschön" ist, dies
funktioniert dann wie folgt:

\begin{enumerate}[label=(\arabic*)]
  \item Nimm an, dass die homogene Lösung $f_h = z_1 f_1 + z_2 f_2$ ist
  \item Versuche nun $f_p = z_1(x) f_1 + z_2(x) f_2$
  \item Löse das folgende System
  \begin{align*}
    z_1'(x) f_1 + z_2'(x) f_2 &= 0\\
    z_1'(x) f_1' + z_2'(x) f_2' &= b(x)
  \end{align*}
  Hierbei gehen wir wie folgt vor, wobei $x_0 \in \X$:
  \begin{align*}
    W(x) &= f_1(x) f_2'(x) - f_2(x) f_1'(x) \neq 0\\
    \Rightarrow z_1' &= \frac{-f_2 b}{W} \; , z_2' = \frac{-f_1 b}{W}\\
    \Rightarrow f_p(x) &= -f_1(x) \int_{x_0}^x \frac{f_2(t) b(t)}{W(t)} dt + f_2(x) \int_{x_0}^x \frac{f_1(t) b(t)}{W(t)} dt
  \end{align*}
\end{enumerate}